
%-----------------------------------------------------------------------------%
\subsection{Simulasi Prediksi menggunakan Studi Kasus}
%-----------------------------------------------------------------------------%
\subsection{Simulasi Prediksi Menggunakan Studi Kasus}

Simulasi prediksi pada bagian ini dilakukan menggunakan data riwayat belajar pengguna yang terdiri dari hasil review (\textit{response history}), interval review (\textit{time history}), probabilitas recall sebelumnya (\textit{p-history}), serta tingkat kesulitan kosakata (\textit{difficulty}). Tujuan simulasi ini adalah menunjukkan bagaimana model GRU-HLR memproses riwayat belajar pengguna untuk menghasilkan prediksi \textit{half-life}, probabilitas recall, dan jadwal review berikutnya.

\subsubsection{Pembentukan Vektor Input}

Riwayat belajar pengguna direpresentasikan sebagai urutan vektor input:

\begin{equation}
x_t=[r_t,t_t,p_t,d_t]
\end{equation}

dengan data sebagai berikut:

\begin{equation}
x_1=[1,0,0.86,0.5]
\end{equation}

\begin{equation}
x_2=[1,1,0.88,0.5]
\end{equation}

\begin{equation}
x_3=[1,3,0.91,0.5]
\end{equation}

\begin{equation}
x_4=[0,7,0.72,0.5]
\end{equation}

Pada penelitian ini, setiap vektor input terdiri atas empat fitur, yaitu hasil review ($r_t$), interval review ($t_t$), probabilitas recall sebelumnya ($p_t$), dan tingkat kesulitan kosakata ($d_t$).

\subsubsection{Perhitungan Update Gate}

Untuk setiap \textit{timestep}, model menghitung \textit{update gate} menggunakan persamaan:

\begin{equation}
z_t=\sigma(W_zx_t+U_zh_{t-1}+b_z)
\end{equation}

Nilai \textit{update gate} menentukan seberapa besar informasi dari \textit{hidden state} sebelumnya dipertahankan untuk membentuk representasi memori baru.

\subsubsection{Perhitungan Reset Gate}

Selanjutnya dihitung \textit{reset gate}:

\begin{equation}
r_t=\sigma(W_rx_t+U_rh_{t-1}+b_r)
\end{equation}

\textit{Reset gate} digunakan untuk menentukan seberapa besar informasi historis yang digunakan kembali dalam proses pembentukan memori baru.

\subsubsection{Perhitungan Candidate Hidden State}

\textit{Candidate hidden state} dihitung menggunakan persamaan:

\begin{equation}
\tilde{h_t}
=
\tanh
\left(
W_hx_t
+
U_h(r_t \odot h_{t-1})
+
b_h
\right)
\end{equation}

Nilai ini merepresentasikan kandidat memori baru yang akan dipertimbangkan oleh model.

\subsubsection{Perhitungan Hidden State}

\textit{Hidden state} akhir pada setiap \textit{timestep} dihitung dengan menggabungkan \textit{hidden state} sebelumnya dan \textit{candidate hidden state}:

\begin{equation}
h_t
=
(1-z_t)\odot h_{t-1}
+
z_t\odot \tilde{h_t}
\end{equation}

Proses ini dilakukan secara berulang untuk seluruh urutan interaksi pengguna:

\begin{equation}
h_0 \rightarrow h_1 \rightarrow h_2 \rightarrow h_3 \rightarrow h_4
\end{equation}

Setelah seluruh data diproses, diperoleh \textit{hidden state} terakhir ($h_4$) yang berisi representasi memori pengguna berdasarkan seluruh riwayat pembelajaran.

\subsubsection{Prediksi Half-Life}

\textit{Hidden state} terakhir kemudian dimasukkan ke \textit{output layer} untuk menghasilkan prediksi \textit{half-life}:

\begin{equation}
\hat{h}
=
W_oh_4+b_o
\end{equation}

Karena target model berupa nilai \textit{half-life} yang harus bernilai positif, hasil transformasi tersebut dipetakan menjadi prediksi akhir:

\begin{equation}
\text{HalfLife}
=
f(\hat{h})
\end{equation}

dengan $f$ merupakan fungsi aktivasi pada \textit{layer} keluaran.

Sebagai ilustrasi, misalkan \textit{hidden state} terakhir yang dihasilkan model adalah:

\begin{equation}
h_4=
\begin{bmatrix}
0.87\\
-0.42\\
0.11\\
\vdots\\
0.55
\end{bmatrix}
\end{equation}

Maka prediksi \textit{half-life} diperoleh melalui:

\begin{equation}
\text{HalfLife}
=
W_oh_4+b_o
=
6.50
\end{equation}

hari.

\subsubsection{Perhitungan Probabilitas Recall}

Setelah nilai \textit{half-life} diperoleh, probabilitas recall dihitung menggunakan persamaan peluruhan memori:

\begin{equation}
R(t)
=
2^{-t/h}
\end{equation}

dengan:

\begin{itemize}
    \item $R(t)$ = probabilitas mengingat
    \item $t$ = jumlah hari sejak review terakhir
    \item $h$ = \textit{half-life} hasil prediksi model
\end{itemize}

Misalnya untuk interval 3 hari dan \textit{half-life} 6,50 hari:

\begin{equation}
R(3)
=
2^{-3/6.5}
\end{equation}

\begin{equation}
R(3)
=
0.726
\end{equation}

Artinya pengguna diperkirakan masih dapat mengingat kosakata tersebut dengan probabilitas sebesar 72,6\%.

\subsubsection{Perhitungan Jadwal Review Berikutnya}

Untuk menentukan waktu review berikutnya digunakan target recall sebesar 90\%.

\begin{equation}
t_{next}
=
\frac
{h\ln(0.9)}
{\ln(0.5)}
\end{equation}

Dengan \textit{half-life} sebesar 6,50 hari:

\begin{equation}
t_{next}
=
\frac
{6.50\times\ln(0.9)}
{\ln(0.5)}
=
0.987
\end{equation}

Hasil tersebut dibulatkan menjadi:

\begin{equation}
t_{next}
=
1
\end{equation}

hari.

Dengan demikian sistem akan merekomendasikan pengguna untuk melakukan review kembali satu hari setelah review terakhir agar probabilitas recall tetap berada di atas target yang ditentukan.

